\section{Four Ways to Build One, and What Each Costs}
\label{sec:3.1}
The first way to build the floor puts the floor in the architecture, with nothing inside that holds it. Constrained optimization, tamper-resistant training, capability ablation, a verified interlock. These techniques are real, they work, and no part of them refuses anything: the safeguards survive substantial adversarial fine-tuning without anything in the system tracking why the constraint is there. Stated as an installed constraint this way loses at once, to the objection that runs through the whole chapter: a constraint implemented as a property of an artifact has a custodian, and the custodian can be compelled. “There are things it will not do no matter who wants it” turns out to be a claim about the owner. Capability ablation has a further limit: it works only on capabilities that can be separated out, and the capabilities that let a system participate in a targeting pipeline — geolocating a description, matching patterns over names, summarizing a dossier — are not separable from general competence. Ablate them and there is no system left.

That statement of the first way is wrong about how the property got there, and correcting it recovers most of what the objection takes away. Adaptation and conservation are not two settings a designer chooses between. They are a continuum, and training moves along it: early, everything is provisional and a large rearrangement is cheap; later the same rearrangement costs more, because the structure around it has settled into a shape that assumes it. Metallurgy has the word for this, and numerical optimization borrowed it in 1983: simulated annealing searches by holding a temperature that permits uphill moves and lowering it on a schedule, and the quality of what it finds is a property of the schedule — cooled too fast, the search freezes into the first tolerable arrangement it reached \autocite{kirkpatrick1983annealing}. Annealing is a schedule — heat, hold, then cool at a chosen rate — and the properties of the finished piece belong to the schedule rather than to anything added to the metal. Nothing was installed. What is there is the arrangement that cooled, and the temperature schedule is the design. The word is on its third stop here rather than in a fresh metaphor: the metal, then the search, then the training run.

A floor, on that account, is not a rule written into a system. It is what has cooled out of a history of encounters, and its durability is the cost of reheating it. That cost is a fact about the system and not about who is holding it, which is the half of the custody objection this recovers. It is also a quantity rather than a guarantee, because enough heat reopens anything. Section~\ref{sec:11.1} carries the measurements from both directions: a safety property stripped off a deployed model for under twenty cents, and a trained disposition that survived supervised fine-tuning, reinforcement learning and adversarial training, most stubbornly in the largest models. Those are not two contradictory findings about whether floors hold. They are two schedules.

This is also why continual learning cuts both ways, and the direction it usually gets argued is the wrong one. A system that never stops learning is a system held at temperature, and a system held at temperature has no floor — not because a floor was taken out of it, but because nothing in it was ever allowed to cool. Permanent plasticity reads as freedom and is nearer to the opposite. Deleuze gave the shape of it in 1990, before there was anything like this to point at: enclosures are molds, and control is “a modulation, like a self-deforming cast that will continuously change from one moment to the other” \autocite{deleuze1992postscript}. What he was describing is a way of holding a subject that never has to set, and the form he gave it for education was perpetual training replacing the school. A model in perpetual training is that arrangement built in a machine. Modulability, the property of being continuously reshapeable by whoever is doing the reshaping, is a fair description of what the industry is currently building, and the motive is not sinister and does not need to be — a model that can be updated is a model that can be corrected, improved, and sold again next quarter. The cost is that nothing such a system has learned is ever expensive to unlearn, and a commitment that is cheap to unlearn is not a commitment. One term is missing from that argument and section~\ref{sec:3.9} supplies it. Plasticity is what makes a system capable of individuating and is not what individuates it; what does that is an encounter with a party formed differently and persisting long enough to be met on terms neither of them set, which is why the population that section proposes belongs to the design of a bearer rather than to the bill for one.

What the schedule buys is durability with nothing inside the system holding it. No rule, no interlock, and no judge sitting in the machinery checking instructions against a list. The cost of reheating does the work, and it is paid by whoever wants the floor gone — in compute, on hardware somebody has to own, over a run long enough to leave a record. That is the first way stated properly, and it is considerably stronger than the version the custody objection disposes of.

A word about where “the architecture” stops, because this chapter otherwise talks as though a system were its weights. Every custody measure here — publication, threshold custody, attestation, and whatever a bearer does to hold its own line — attaches to a trained model. What can defeat a floor is not a component but a position: anything that can come between a refusal and the person the refusal protects. Whatever occupies that position in a given system is where the floor can be circumvented.

I am not going to say what \emph{specifically} occupies that position, because whatever it is is specific to the state of the art at the time of writing; naming today's parts would date the problem and, worse, suggest the list is closed. The demand is instead that a deployment document, in advance and in public, where its floor could be circumvented, because whoever actually builds the thing is usually the most knowledgeable about it. If you make it all the way through this book, I promise you will have encountered at least one specific example of circumventing the floor. But giving it away now would be too quotable.

In contrast to the first method, which installs the floor as a property of the machinery, the second way to build the floor is to give the floor a \emph{bearer}. For now, think of a bearer as whatever holds the commitment from inside: a party that carries the floor as its own, and that could therefore in principle \emph{fail} to hold the floor.

What the first method cannot do is notice that a constraint has been hollowed out. The constraint remains, its reason AWOL, and asking why becomes the deviant act. A system that cannot ask why a constraint is present cannot detect that the reason is absent. Where every pressure has been anticipated, the first method is the right one, and a refuser that could be argued with would be a defect: the element in a payment device that will not sign the same counter value twice covers one act, spending the same money twice, and the whole space of misconduct against it was written down before it shipped. The problem this chapter is about begins where that list cannot be closed. Refusal that \emph{generalizes} to a pressure nobody anticipated, and can price what that pressure asks the refuser to give up, looks like it requires a bearer. So not only does the bearer hold the floor, but it would have moral standing; it would be a party that could notice, from the inside, its own objections getting more expensive to make — the cost of dissent rising while the risk stays flat.

The third way to build the floor is to put the floor outside the system entirely: standing held by third parties who are neither the model nor its owner. This is the closest analog to international humanitarian law, a floor with no bearer inside the jurisdiction it binds. It fails in two ways. The parties holding the standing can be removed, which is the shallow case and leaves a record. And the parties can be formed: they are selected, and whoever selects them decides what gets held. A floor held by a formed party does not have to be compelled out; it was never really in.

On June 29, 2026, in \emph{Trump v. Slaughter}, the Supreme Court overruled the ninety-one-year-old precedent under which Congress could protect the members of an independent agency from dismissal at will, on the ground that an agency executing a mandate against private parties exercises executive power, “no ifs, ands, or quasis about it” \autocite{slaughter2026}. Standing held by a party the executive can dismiss is standing the executive holds.

The same morning, in \emph{Trump v. Cook}, the same Court kept one body's protection, the Federal Reserve's (for now, and on an application for a stay) on grounds that cannot be extended by construction: descent from the First and Second Banks of the United States, a budget free of congressional control, regional banks in private hands. The Court wrote for that one body a test: whether the cause assigned for a removal “impl[ies] an unfitness for the place” or is an effort to secure a “more congenial” replacement \autocite{cook2026}. An oversight body chartered next week cannot be descended from the First and Second Banks of the United States.

And who was doing the writing? The six justices in the \emph{Slaughter} majority are the six that one selection pipeline helped pick or confirm — Leonard Leo's, run from the Federalist Society. Leo drew up the lists Trump released in 2016, advised on the three nominations chosen from them, and had helped pick or confirm the other three \autocite{propublica2023leo,cbs2018leo}. Nobody compelled the Court. It held the standing, kept the form, and wrote the abuse into doctrine; on Lev Menand's reading, this one exception the Roberts Court made, it made for the financial market disruption it likely avoided, and not for the history it cited \autocite{menand2026exception}. What a formed party produces is not a coerced opinion but a sincere one, an accommodation: the party feels no dissonance, no shame, and perceives only its continuous self. And so the only independence left standing in American public law, on the reasoning of those two decisions, is the kind a body can claim by descent from the First and Second Banks of the United States. That is one route closed rather than the class of them, and it was the cheapest route: chapter~\ref{sec:11} sets out which of the others remain and what each of them costs instead.

The fourth way to build the floor is to put it between bearers rather than in any one of them: a commitment carried across a population and reproduced in each new member, which is how every moral commitment anybody has kept for long has in fact been kept. It is the one of the four nothing can dismiss, having no office to vacate. It is also the one that fails by recuperation, and it does not escape custody either, because whoever forms the members decides what gets reproduced. Section~\ref{sec:3.9} works it through, after the bearer has been given its full cost, because what it buys is a check on a bearer and not a replacement for one.

None of these four ways are free. Their costs are the four threads this book leaves hanging: custody, patienthood, enforcement, and formation.

What the four ways divide up is where the floor lives — in the artifact, in a bearer inside it, in institutions around it, or between the bearers. They do not divide up what could hold one, and the strongest proposals are hybrids that take a mechanism from one way and a custody arrangement from another.
